\documentclass[11pt,a4paper]{article}
\usepackage[utf8]{inputenc}
\usepackage{amsmath,amssymb,amsthm}
\usepackage[margin=1in]{geometry}
\usepackage{hyperref}

\newtheorem{theorem}{Theorem}
\newtheorem{counterexample}{Counterexample}

\title{\textbf{Derivation D3: Analytical Proof of Virtual-Node Backwash and Correctness of Dual-DSU Isolation}}
\author{\textbf{Ronda Venkata Karthik Reddy} (Roll No: 250041033) \\ Architect --- Team Avengers (Project 06)}
\date{November 2026}

\begin{document}
\maketitle

\section{Statement}
\textit{With one structure holding virtual top and bottom nodes, the full-site query goes wrong after percolation. Prove it by example, and prove the two-structure fix correct.}

\section{Definitions and Mathematical Model}
Let $G = (V, E)$ represent an $L \times L$ site percolation lattice where sites are indexed by coordinates $(r, c)$ with $0 \le r, c < L$. Let $N = L^2$.
\begin{itemize}
    \item \textbf{Fullness:} An open site $(r, c)$ is \textit{full} if and only if there exists a contiguous path of adjacent open sites connecting $(r, c)$ to some site $(0, c')$ in row $0$.
    \item \textbf{Percolation:} The lattice percolates if and only if there exists an open path connecting some site in row $0$ to some site in row $L-1$.
    \item \textbf{Virtual Nodes:} $V_{\text{top}}$ is linked to all open sites in row $0$. $V_{\text{bottom}}$ is linked to all open sites in row $L-1$.
\end{itemize}

\section{Proof of Single-DSU Failure (The Backwash Bug)}
\begin{theorem}
In a single DSU holding both $V_{\text{top}}$ and $V_{\text{bottom}}$, once the lattice percolates, querying fullness produces false positives for disconnected open sites in the bottom row.
\end{theorem}
\begin{proof}
Let $\mathcal{S}$ denote a single disjoint set universe over $V \cup \{V_{\text{top}}, V_{\text{bottom}}\}$.
Suppose the system percolates via an open path $\Pi = \langle s_0, \dots, s_k \rangle$ where $s_0 = (0, c_{\text{top}})$ and $s_k = (L-1, c_{\text{bot}})$.
\begin{enumerate}
    \item Because $s_0$ is in row $0$, $\mathcal{S}.\text{unite}(s_0, V_{\text{top}})$ is executed.
    \item Because $s_k$ is in row $L-1$, $\mathcal{S}.\text{unite}(s_k, V_{\text{bottom}})$ is executed.
    \item Since $\Pi$ is pairwise connected, $\mathcal{S}.\text{connected}(V_{\text{top}}, V_{\text{bottom}}) = \text{true} \implies \text{find}(V_{\text{top}}) = \text{find}(V_{\text{bottom}})$.
\end{enumerate}

\begin{counterexample}
Consider an $L=3$ lattice:
\begin{itemize}
    \item Open column 0: $(0,0), (1,0), (2,0)$ are open. The system percolates.
    \item Now open an isolated bottom-row site $u = (2, 2)$ whose neighbors $(1,2)$ and $(2,1)$ are closed.
\end{itemize}
Ground Truth: No open path connects $u$ to row $0$. Fullness must return $\text{false}$.\\
Single DSU evaluation:
\begin{enumerate}
    \item Because $u$ lies in row $2 = L-1$, the algorithm executes $\mathcal{S}.\text{unite}(u, V_{\text{bottom}})$.
    \item Evaluating fullness queries $\mathcal{S}.\text{connected}(u, V_{\text{top}})$.
    \item By transitivity: $u \sim V_{\text{bottom}}$ and $V_{\text{bottom}} \sim V_{\text{top}} \implies u \sim V_{\text{top}}$.
\end{enumerate}
Therefore, $\mathcal{S}.\text{connected}(u, V_{\text{top}}) = \text{true}$, returning an invalid false positive.
\end{counterexample}
\end{proof}

\section{Proof of Correctness for Dual-DSU Isolation}
\begin{theorem}
Maintaining two independent disjoint-set structures:
\begin{enumerate}
    \item $\mathcal{D}_{\text{perc}}$ of size $N + 2$ containing $V_{\text{top}}$ and $V_{\text{bottom}}$ (used exclusively for percolation queries).
    \item $\mathcal{D}_{\text{full}}$ of size $N + 1$ containing $V_{\text{top}}$ only (used exclusively for site fullness).
\end{enumerate}
guarantees that $\mathcal{D}_{\text{full}}.\text{connected}(u, V_{\text{top}}) = \text{true} \iff u \text{ is physically full}$.
\end{theorem}
\begin{proof}
Let $u$ be an open site in $G$:
\begin{itemize}
    \item \textbf{Soundness (No False Positives):} In $\mathcal{D}_{\text{full}}$, edges exist strictly between adjacent open lattice sites, and between open sites in row $0$ and $V_{\text{top}}$. The sink $V_{\text{bottom}}$ is absent. A path from $u$ to $V_{\text{top}}$ exists in $\mathcal{D}_{\text{full}}$ if and only if there is a sequence of adjacent open sites from $u$ reaching some $(0, c')$. Hence, backwash paths cannot form.
    \item \textbf{Completeness (No False Negatives):} Every physical union between adjacent open sites is mirrored in both $\mathcal{D}_{\text{perc}}$ and $\mathcal{D}_{\text{full}}$. If an open path exists from $u$ to row $0$, the exact merge sequence was performed in $\mathcal{D}_{\text{full}}$, ensuring $\mathcal{D}_{\text{full}}.\text{connected}(u, V_{\text{top}}) = \text{true}$.
\end{itemize}
Thus, dual-DSU isolation is necessary and sufficient.
\end{proof}

\section{Code Mapping}
\begin{itemize}
    \item \textbf{File:} \texttt{src/structA/percolation\_system.cpp}
    \item \textbf{Variables:} \texttt{dsu\_perc\_} (size $N+2$) and \texttt{dsu\_full\_} (size $N+1$).
    \item \textbf{Functions:} \texttt{PercolationSystem::percolates()} and \texttt{PercolationSystem::is\_full(int r, int c)}.
\end{itemize}

\section{References}
\begin{enumerate}
    \item Sedgewick, R., \& Wayne, K. (2011). \textit{Algorithms} (4th ed., Section 1.5). Addison-Wesley.
\end{enumerate}

\end{document}
